\documentclass[11pt]{scrartcl}

\usepackage[sexy]{evan}
\usepackage{import}
\usepackage[table]{xcolor}
\usepackage{xfp}

\newcommand{\gad}{\textcolor{yellow}{$\bigstar$}}
\newcommand{\bad}{\textcolor{red}{$\bigstar$}}

\definecolor{dcol0}{HTML}{C8E6C9}
\definecolor{dcol1}{HTML}{D4E9B3}
\definecolor{dcol2}{HTML}{E5ED9A}
\definecolor{dcol3}{HTML}{FFF59D}
\definecolor{dcol4}{HTML}{FFE082}
\definecolor{dcol5}{HTML}{FFCC80}
\definecolor{dcol6}{HTML}{FFAB91}
\definecolor{dcol7}{HTML}{F49890}
\definecolor{dcol8}{HTML}{E57373}
\definecolor{dcol9}{HTML}{D32F2F}

\makeatletter
\newcommand{\getcolorname}[1]{dcol#1}
\makeatother

\newcommand{\dif}[1]{%
    \edef\colorindex{\number\fpeval{floor(#1)}}%
    \edef\fulltext{#1}%
    \colorbox{\getcolorname{\colorindex}}{%
        \ifnum\colorindex>8
            \textbf{\textcolor{white}{\,\fulltext\,}}%
        \else
            \textbf{\textcolor{black}{\,\fulltext\,}}%
        \fi
    }%
}
% Variable para dificultad (inicial 0)
\newcommand{\thmdifficulty}{0}

% Comando para asignar dificultad antes del problema
\newcommand{\problemdiff}[1]{\renewcommand{\thmdifficulty}{#1}}

% Estilo del problema que incluye dificultad antes del título
\declaretheoremstyle[
    headfont=\color{blue!40!black}\normalfont\bfseries,
    headformat={%
      \dif{\thmdifficulty}\quad \NAME~\NUMBER\ifx\relax\EMPTY\relax\else\ \NOTE\fi
    },
    postheadspace=1em,
    spaceabove=8pt,
    spacebelow=8pt,
    bodyfont=\normalfont
]{problemstyle}

    \declaretheorem[style=problemstyle,name=Problema,sibling=theorem]{problema}
    \declaretheorem[style=problemstyle,name=Problema,numbered=no]{problema*}

\title{Alg Manip}
\subtitle{OMMJAL}
\author{Emmanuel Buenrostro}

\begin{document}

\maketitle.

\section{Tecnicas}

\subsection{SFFT}

Simon's Favorite Factoring Trick. Este truco sirve para factorizas cosas de la forma: (Esto lo vamos a nombrar SFFT con $(x,y,n,m)$)
\[ xy+xn+ym \]

Para esto vamos a usar que 
\[ xy+xn+ym+nm =(x+m)(y+n)\]
Es bastante parecido, entonces el truco es agregarle $nm$ para poder factorizarlo, entonces

\[ xy+xn+ym = (x+m)(y+n)-nm \]

Algunos casos especiales muy comunes son:

\[xy+x+y=(x+1)(y+1)-1  \text{ y } xy-x-y=(x-1)(y-1)-1\]

Si tienes cosas de la forma
\[ kxy+ax+by=0 \]
Tambien puedes aplicarlo, pero de una forma más curiosa que es multiplicando por $k$. 

\[ k^2xy +akx +bky =0\]
y usando SFFT con $(kx,ky,+a, +b)$.
\[(kx+b)(ky+a)=ab\]

\begin{example}
    [Brilliant]
    Encuentra las soluciones enteras a la ecuación $xy+9(x+y)=86$
\end{example}

\begin{example}
    [Brilliant]
    Encuentra las soluciones enteras a la ecuación $2b+3c=5bc$.
\end{example}



\newpage
\section{Problemas}


\begin{problem} %OMM 2023/1
    Encuentra los enteros positivos de cuatro digitos, tal que la suma del cuadrado
    de los digitos sea igual a el doble de la suma de los digitos.
\end{problem}




\begin{problem}
    %OMM 2010/1
    Encuentra todas las ternas de números naturales $(a,b,c)$ tales que $abc=a+b+c+1$.
    
\end{problem}

\begin{problem}
    %Lista ommbc
    Resuelve el sistema de ecuaciones:
    \begin{align*}
        x^3-y^3&=19(x-y)\\
        x^3+y^3&=7(x+y)
    \end{align*}
\end{problem}


\begin{problem}
Sean $p,q$ dos números primos. Resuelve en enteros la siguiente ecuación:
\[ \frac 1x +\frac 1y = \frac{1}{pq}\]
\end{problem}

\begin{problem} %sophie-germain
    Factoriza $a^4+4b^4$
\end{problem}

\begin{problem}
    %lista ommbc
    Demuestra que si $n$ es la suma de dos cuadrados, entonces $2m$ tambien lo es.
\end{problem}

(Si no has visto log saltate este problema)
\begin{problem}
    %Girls Math at Yale
    Sean $x,y$ reales positivos tal que $\log_2 x =\log_x y = \log_y 256$. ¿Cuánto vale $xy$?
\end{problem}   

\begin{problem} %Lista ommbc
    Resolver el sistema de ecuaciones:
    \begin{align*}
        x^2+xy+y^2&=4 \\
        x+xy+y&=2 
    \end{align*}
\end{problem}

\begin{problem}
    %lista ommbc
    Encuentra las soluciones reales a:
    \begin{align*}
        x^3+y^3 &=1 \\
        x^2y+2xy^2+y^3 &= 2 
    \end{align*}    
\end{problem}

\begin{problem}
    %OMM 2022/1
    Un número $x$ es Tlahiuca si existen primos positivos distintos $p_1,p_2,\dots,p_k$ tales que 
\[x=\frac{1}{p_1}+\frac{1}{p_2}+\cdots+\frac{1}{p_k}.\]
Deterimina el mayor número Tlahuica $x$ que satisface las dos siguientes propiedades:
 \begin{itemize} 
 \item  $0 < x < 1$. 
 \item  Existe un número entero $0< m\leq 2022$ tal que $mx$ es un número entero. 
 \end{itemize} 
    
\end{problem}




\begin{problem}
    %IMOSL 2013 A1
    Sea $n$ un entero positivo, y sean $a_1, \ldots, a_{n-1}$ números reales arbitrarios.
    Las sucesiones $u_0, u_1, \ldots , u_n$ y $v_0, v_1, \ldots, v_n$ como $u_0=u_1=v_0=v_1=1$ y 
    \[u_{k+1}=u_k+a_ku_{k-1} \text{ y } v_{k+1}=v_k+a_{n-k}v_{k-1}\]
    para $k=1,2,\ldots, n-1$. \\
    Prueba que $u_n=v_n$.
\end{problem}


\end{document}